\documentclass[10pt,a4paper]{amsart}
\usepackage[margin=19mm]{geometry}
\usepackage{amsmath,amssymb,mathtools}
\usepackage[scr=rsfs,cal=euler]{mathalfa}
\usepackage{xfrac}
\usepackage[unicode,hidelinks,bookmarks=false]{hyperref}
\newcommand{\ie}{i.e.\ }
\newcommand{\cf}{cf.\ }
\newcommand{\resp}{respectively\ }
\newcommand{\Subsection}{\subsection}
\providecommand{\nobreakdash}{\nobreak\hbox{-}}
\setlength{\emergencystretch}{2em}
\multlinegap0pt
\usepackage{bm}
\usepackage{bbm}
\usepackage{dsfont}
\usepackage{xspace}
\usepackage{cancel}
\usepackage{mathtools}
\usepackage{amsthm}
\makeatletter
\@ifundefined{theorem}{\newtheorem{theorem}{Theorem}}{}
\@ifundefined{conjecture}{\newtheorem{conjecture}[theorem]{Conjecture}}{}
\@ifundefined{proposition}{\newtheorem{proposition}[theorem]{Proposition}}{}
\makeatother
\newcommand{\Irr}{{\mathrm {Irr}}}
\newcommand{\QQ}{{\mathbb Q}}
\title[Problems on the conductor of finite group characters]{Problems on the conductor of finite group characters}
\author{Nguyen N. Hung}
\date{Source: arXiv:2512.18042v3 (CC BY 4.0)}
\begin{document}
\maketitle

\noindent\emph{Source and licence.} Problems on the conductor of finite group characters. Version arXiv:2512.18042v3 (CC BY 4.0), distributed under the Creative Commons Attribution 4.0 International licence, as recorded in the arXiv licence metadata for this exact version and shown on the abstract page https://arxiv.org/abs/2512.18042v3. Full rights notice: licenses/arxiv-2512.18042-CC-BY-4.0.txt.\par\smallskip

\noindent\textit{Editorial scope.} Counted unit: the Proposition whose source label is \texttt{prop:1} in the article (source0.tex lines 550--553, source label \texttt{prop:1}), delivered together with the article's own complete original proof of it (source0.tex lines 555--593, one \texttt{proof} environment of 39 lines). No mathematics is written, repaired or re-derived: statement and proof are the article's own text line for line. Required context retained uncounted: conj:stronger (lines 541--548). Every definition, display and label of those lines is retained unchanged. Omitted from this package and not redistributed: 1--540, 549--549, 554--554, 594--1539. Nothing inside a retained excerpt is omitted or altered: every retained body is delivered exactly as the article's lines are written, so no omission receipt of any kind accompanies this package. Citations: the retained excerpts carry no citation command at all, so no bibliography entry of the article is transcribed and no reference is redistributed. Reference adaptation: none was needed and none was made; every reference inside the delivered unit resolves inside it, so no unresolved reference or citation is produced. Printed slips kept literal: no printed word, symbol, hypothesis or number of the counted statement or of its proof was altered. Layout and class: the article's own class is \texttt{amsart} with packages including latexsym, enumerate, amssymb, xcolor, mathrsfs, amsmath, amsthm, amsfonts, mathabx, hyperref, inputenc, arydshln; the wrapper is the lane's pinned \texttt{amsart} editorial wrapper at 10pt on A4, which restates the theorem-like environments the retained text needs and adds the article's own macro definitions that the retained text needs (\texttt{\textbackslash Irr}, \texttt{\textbackslash QQ}), with extra wrapper packages (bm, bbm, dsfont, xspace, cancel, mathtools). The delivered numbering is the unit's own. Assets: no retained span contains a figure environment, a graphics-inclusion command, a PDF-page inclusion, a table or a TikZ picture; no image file is copied into this package, the package declares no asset dependency and no image file is redistributed with it. Licence: version arXiv:2512.18042v3 of this article is distributed under the Creative Commons Attribution 4.0 International licence, as recorded in the arXiv licence metadata for this exact version and shown on the abstract page https://arxiv.org/abs/2512.18042v3. A full rights notice is delivered at licenses/arxiv-2512.18042-CC-BY-4.0.txt. Characters: every retained excerpt is pure ASCII, so the delivered engine, XeLaTeX with its default Latin Modern text font, reports no missing character.
\begin{conjecture}\label{conj:stronger}
Let $G$ be a finite group and $\chi\in\Irr(G)$ be non-linear. Then
there exist a proper subgroup $H<G$ and $\psi\in\Irr(H)$ such that
\[
\QQ(\chi) \subseteq \QQ(\psi) \quad \text{and} \quad [\QQ(\psi) :
\QQ(\chi)] \le \frac{\chi(1)}{\psi(1)}.
\]
\end{conjecture}

\begin{proposition}\label{prop:1}
Conjecture~\ref{conj:stronger} is equivalent to the combination of
the cyclotomic deficiency conjecture and Feit's conjecture.
\end{proposition}

\begin{proof}
We first show that Conjecture~\ref{conj:stronger} implies the
cyclotomic deficiency conjecture. We proceed by induction on $|G|$.
Note that the cyclotomic deficiency conjecture is trivial for linear
characters, so we may assume that $\chi(1) > 1$. By hypothesis,
there exists $H < G$ and $\psi \in \Irr(H)$ such that $\QQ(\chi)
\subseteq \QQ(\psi)$ and $[\QQ(\psi) : \QQ(\chi)] \le
\chi(1)/\psi(1)$. It follows that
\[
[\QQ_{c(\chi)} : \QQ(\chi)] \le [\QQ_{c(\psi)} : \QQ(\chi)] =
[\QQ_{c(\psi)} : \QQ(\psi)] \, [\QQ(\psi) : \QQ(\chi)].
\]
By the induction hypothesis, we have $[\QQ_{c(\psi)} : \QQ(\psi)]
\le \psi(1)$. Hence,
\[
[\QQ_{c(\chi)} : \QQ(\chi)] \le \psi(1) \cdot
\frac{\chi(1)}{\psi(1)} = \chi(1),
\]
as desired.

To see that Conjecture~\ref{conj:stronger} also implies Feit's
conjecture, we apply Conjecture~\ref{conj:stronger} repeatedly to
obtain a linear character $\lambda$ of some subgroup $K$ of $G$ such
that
\[
\mathbb{Q}(\chi) \subseteq \mathbb{Q}(\lambda) =
\mathbb{Q}_{o(\lambda)}.
\]
Hence $c(\chi)$ divides $o(\lambda)$. Since $o(\lambda)$ equals the
order of some element of $K/K'$, it follows that $c(\chi)$ divides
the order of some element $x \in K$. Moreover, some power of $x$
will have order $c(\chi)$.


For the converse implication, let $g \in G$ be an element of order
$|g| = c(\chi) =: a$. If $G = \langle g \rangle$, then there is
nothing to prove. Otherwise, we may take $H = \langle g \rangle$ and
let $\psi$ be an irreducible character of $H$ of order $a$.
\end{proof}

\end{document}
